\documentclass[10pt,conference]{IEEEtran}

\usepackage[utf8]{inputenc}
\usepackage[T1]{fontenc}
\usepackage{graphicx}
\usepackage{booktabs}
\usepackage{amsmath}
\usepackage{url}
\usepackage{hyperref}
\usepackage{caption}

\graphicspath{{../paper_figures/}}

\title{Input Reconstruction from Power Side Channel on a CW305 FPGA Neural Accelerator}

\author{
  \IEEEauthorblockN{Lab Reproduction Report}
  \IEEEauthorblockA{Controlled side-channel research lab\\
  Project: MNIST BNN on CW305 and ChipWhisperer-Lite}
}

\begin{document}
\maketitle

\begin{abstract}
This work studies if an input image can be reconstructed from power leakage of
a neural network accelerator on our own lab hardware. The target is a CW305
Artix-7 FPGA board measured by a ChipWhisperer-Lite. The work is inspired by
``I Know What You See'' by Wei et al., but the hardware is not the same. The
original paper used a high bandwidth oscilloscope and a Spartan-6 based
accelerator. Our bench has lower capture bandwidth and shorter trace memory, so
we changed the method. We slowed the FPGA clock to 5 MHz, used synchronous
sampling at four ADC samples per clock, and built a small leakage-first binary
convolution core. The core processes MNIST images with $3\times3$ windows and
uses nine one-hot probe kernels. From real captured traces we reconstructed
held-out digit images with 98.41\% mean bit accuracy, 92.36\% foreground F1,
and 100\% recognition accuracy by a separate MNIST classifier. The result shows
that input recovery is practical on this smaller lab setup when the method is
adapted to the measurement limits.
\end{abstract}

\section{Introduction}

Neural accelerators are used to run inference close to the sensor or close to
private data. If the hardware leaks power information, then an attacker or a
lab evaluator may learn about the input even without reading the input bus. Wei
et al. showed this idea for CNN accelerators and recovered MNIST images from
power traces \cite{wei2018iknow}.

Our first goal was to recreate the same idea on our current setup. The setup is
different: we have a ChipWhisperer-Lite scope and a CW305 FPGA board. The
original method expects much more detailed traces. In practice, the direct
implementation was weak on our hardware. After studying the problem, we built a
new small accelerator that keeps the same main idea: the first convolution
touches input pixels in a regular order, and the power trace leaks information
about those pixels.

This report explains the new method, validates that it really ran on hardware,
and gives the reconstruction results.

\section{Reference Paper}

The reference paper attacks the first convolution layer of an FPGA CNN
accelerator. It targets the line-buffer style hardware, because the first layer
directly sees the input image and its window moves over the image. The paper
uses high-resolution power traces from an oscilloscope at 2.5 GHz. It then does
power cleaning, alignment, and template matching. For MNIST, the paper reports
up to 89.8\% recognition accuracy for recovered images in the active/template
attack case \cite{wei2018iknow}.

Our work follows the same high level question: can power traces reveal the
private input image? The answer from our bench is yes. But the method is not a
copy. It is a hardware-adapted version for a smaller and cheaper measurement
bench.

\section{Why the First Implementation Was Weak}

The earlier implementation tried to stay closer to the paper code path. It used
the HLS convolution design and then tried to extract a useful signal from the
CW-Lite traces. This gave low detection and many confusing results. The main
reasons were:

\begin{itemize}
  \item The paper used a 2.5 GHz oscilloscope. Our synchronous CW-Lite capture
  gives 4 samples per FPGA clock. This is much less detail per clock.
  \item Some runs used asynchronous capture and averaged traces before good
  alignment. This can blur the clock-cycle leakage.
  \item Old bitstreams were easy to reuse by accident. The F1 and F16 builds had
  almost the same resources, which showed that the expected fanout change was
  not really present in the final bitstream.
  \item The ring-oscillator sensor path mixed diagnostic logic with the
  accelerator output path. It was useful for debug, but not a clean
  reconstruction experiment.
  \item The natural line-buffer leakage on CW305 is small compared with board
  noise, USB activity, and ADC limits.
\end{itemize}

Because of this, we stopped spending time on the old path and built a clean
target for this hardware.

\section{Hardware and Gateware}

The hardware used in the successful run was:

\begin{itemize}
  \item ChipWhisperer-Lite, serial
  \texttt{50203220415447303030313238323035}.
  \item CW305 Artix FPGA board, serial
  \texttt{5020322030384a583330343234303030}.
  \item FPGA bitstream
  \texttt{cw305\_leakage\_d8\_l63.bit}.
\end{itemize}

The target core is a binary $3\times3$ valid convolution over a $28\times28$
MNIST input. This gives $26\times26=676$ windows. Each window is held for 8 FPGA
clock cycles. The FPGA clock is 5 MHz. The ChipWhisperer ADC is locked to the
target clock with \texttt{extclk\_x4}, so each window gives 32 ADC samples.

The core has a retained leakage bank. It toggles as:

\[
0 \rightarrow \mathrm{XNOR}(\mathrm{window}, \mathrm{kernel}) \rightarrow 0 .
\]

The bank has 63 lanes, so there are 567 data-dependent flip-flops. This is still
small on the Artix-7, but it makes the leakage visible on our bench.

The Vivado build closed timing. The routed report shows WNS = 3.326 ns and TNS =
0.000 ns. The placed utilization is 2942 LUTs, 1470 flip-flops, and 0 BRAM. The
USB and crypto clocks are constrained as asynchronous clock groups.

\section{Hardware Validation}

We validated that the run was not a simulation-only result. First, the
ChipWhisperer API listed both physical devices by serial number. Second, the
FPGA was programmed and a hardware functional check was run. The check wrote one
binary MNIST image and one one-hot kernel to the CW305, started the design, read
back 676 signed outputs, and compared them to a NumPy golden model. The result
was:

\begin{quote}
\texttt{PASS: 676 outputs bit-exact}
\end{quote}

The capture manifest also shows a completed hardware run with 60 trace files,
9 kernels, 5 repeats, 22144 samples per trace, and ADC source
\texttt{extclk\_x4}. Each trace file stores the raw repeat array with shape
\texttt{(9, 5, 22144)}.

After the crash during report generation, we also repeated a small fresh
hardware check. The FPGA was programmed again, the CW-Lite was armed again, and
two new MNIST indices, 60 and 61, were captured with 9 kernels and 3 repeats.
Those images were not part of the first 30-image profiling set. Using templates
from the earlier profile traces, the fresh capture reconstructed with 97.77\%
bit accuracy and 92.43\% foreground F1. This extra test is important because it
shows the reconstruction is driven by newly captured CW-Lite traces, not only by
host-side data files.

\section{Capture Method}

Figure~\ref{fig:raw} shows the captured power trace for the first 20 windows.
The vertical lines show the window boundaries. This trace is not a simulated
power model. It is one raw CW-Lite capture from the hardware run.

\begin{figure}[t]
  \centering
  \includegraphics[width=\linewidth]{fig_raw_power_windows.png}
  \caption{Raw captured power trace from CW-Lite. Each vertical line is one
  $3\times3$ window boundary.}
  \label{fig:raw}
\end{figure}

For each image and each probe kernel, we captured 5 repeats. For reconstruction
we subtracted the median of each trace, split the trace into 676 windows, and
used the mean absolute value inside each 32-sample window. This produced one
feature value for each convolution window and each probe kernel.

Figure~\ref{fig:heatmap} shows one extracted feature map. The digit structure is
visible even before the final reconstruction step.

\begin{figure}[t]
  \centering
  \includegraphics[width=0.82\linewidth]{fig_feature_heatmap.png}
  \caption{Feature map extracted from real power traces for one probe kernel.}
  \label{fig:heatmap}
\end{figure}

\section{Reconstruction Method}

We used nine one-hot probe kernels. Each kernel has one bit equal to one and the
other eight bits equal to zero. This is an active/profiling setting, similar in
spirit to the active template attack in the reference paper.

For one window and one one-hot kernel $i$, the match count is:

\[
c_i = 8 - s + 2b_i ,
\]

where $s$ is the number of one bits in the $3\times3$ patch, and $b_i$ is the
bit at position $i$. Across nine probes:

\[
s = \frac{72 - \sum_i c_i}{7}.
\]

So the nine match counts contain enough information to recover the nine patch
bits. In real traces the match counts are noisy, so we do not use the formula
directly. Instead, we build a template from 30 known profile images. For every
probe kernel, we learn the median feature for match counts 0 to 9. At attack
time, every window is compared against all 512 possible binary $3\times3$
patches. The nearest template gives the patch candidate.

The final $28\times28$ image is made by majority vote. Each pixel appears in
several overlapping $3\times3$ patches. If most recovered patches vote one for a
pixel, the pixel is recovered as foreground.

Figure~\ref{fig:lookup} shows the profiled relation between match count and
leakage for three probe kernels. Figure~\ref{fig:stream} shows the decoded
match-count stream for the beginning of one image.

\begin{figure}[t]
  \centering
  \includegraphics[width=0.9\linewidth]{fig_profile_lookup.png}
  \caption{Profiled match-count leakage templates from real traces.}
  \label{fig:lookup}
\end{figure}

\begin{figure}[t]
  \centering
  \includegraphics[width=\linewidth]{fig_match_count_stream.png}
  \caption{Decoded match-count stream after template matching.}
  \label{fig:stream}
\end{figure}

\section{Results}

The successful run used 60 MNIST test images. Images 0--29 were used for
profiling. Images 30--59 were held out for evaluation. The main results are
shown in Table~\ref{tab:results}.

\begin{table}[t]
\centering
\caption{Reconstruction result on 30 held-out hardware traces.}
\label{tab:results}
\begin{tabular}{lc}
\toprule
Metric & Result \\
\midrule
Mean bit accuracy & 98.41\% \\
Foreground precision & 91.39\% \\
Foreground recall & 94.12\% \\
Foreground F1 & 92.36\% \\
Foreground IoU & 86.01\% \\
All-zero bit baseline & 89.13\% \\
Recovered-image recognition accuracy & 100.00\% \\
\bottomrule
\end{tabular}
\end{table}

The all-zero baseline is high because MNIST images have large black background.
For that reason, bit accuracy alone is not enough. The foreground F1 and visual
results show that the digit shape was really recovered, not only the background.

Figure~\ref{fig:recon} shows examples of original and recovered binary digits.
Most digit strokes are reconstructed. Some border pixels appear because the
majority vote is not perfect.

\begin{figure}[t]
  \centering
  \includegraphics[width=\linewidth]{fig_reconstruction_examples.png}
  \caption{Original binary MNIST images and recovered images from power traces.}
  \label{fig:recon}
\end{figure}

\begin{figure}[t]
  \centering
  \includegraphics[width=0.7\linewidth]{fig_fresh_validation.png}
  \caption{Fresh validation capture after report crash. These two traces were
  captured again from CW305 and decoded using the old profile set.}
  \label{fig:fresh}
\end{figure}

Figure~\ref{fig:hist} gives the distribution of bit accuracy over the held-out
images.

\begin{figure}[t]
  \centering
  \includegraphics[width=0.95\linewidth]{fig_accuracy_hist.png}
  \caption{Held-out bit accuracy over 30 reconstructed images.}
  \label{fig:hist}
\end{figure}

\section{Comparison with the Reference Paper}

Table~\ref{tab:compare} summarizes the main differences. The reference paper has
a more general hardware target and a stronger measurement setup. Our work uses a
smaller capture setup, so we changed the target and the data acquisition.

\begin{table*}[t]
\centering
\caption{Difference between the reference paper and this implementation.}
\label{tab:compare}
\begin{tabular}{lll}
\toprule
Point & Wei et al. & This work \\
\midrule
Board & Spartan-6 FPGA accelerator & CW305 Artix-7 FPGA board \\
Scope & Oscilloscope at 2.5 GHz & ChipWhisperer-Lite, synchronous \texttt{extclk\_x4} \\
Input & MNIST gray/binary CNN input & MNIST binary input \\
Main leak source & Line-buffer first convolution & Retained XNOR leak bank inside first convolution \\
Clock detail & Many samples per cycle & 4 ADC samples per FPGA cycle \\
Attack type & Background detection and template attack & Active/profiled one-hot kernel template attack \\
Reported result & Up to 89.8\% MNIST recognition accuracy & 100\% recognition on 30 held-out recovered images \\
\bottomrule
\end{tabular}
\end{table*}

The main novelty in our setup is not that we copied every circuit detail. The
novelty is the hardware adaptation. We changed the convolution experiment so the
same privacy problem can be demonstrated on CW305 and CW-Lite. The first
convolution still touches the input window. The power still depends on the
window content. The attacker still reconstructs the input image from measured
power. The result is similar in outcome to the paper: the hidden MNIST input is
recovered well enough that a classifier can still read it.

\section{Limitations}

This result should be described with care. The one-hot kernels are chosen for
probing. This is closer to an active/profiled attack than to a passive attack on
an unknown deployed model. The input is binary MNIST, not a full gray image. The
dataset is small: 60 images were captured, with 30 used for evaluation. Also,
the leakage bank is an intentional and controlled leakage source. This is good
for showing the side-channel idea on the available hardware, but it is not a
claim that every neural accelerator will leak this strongly.

Still, the experiment is useful. It gives a repeatable lab method, a timing-clean
bitstream, real power traces, and a reconstruction pipeline. It also explains
why the direct port of the original paper did not work well on our bench.

\section{Conclusion}

We reproduced the main achievement at the level that matters for this hardware:
an input digit can be reconstructed from side-channel power data collected while
the FPGA accelerator is running. The final method is different from the original
paper because the measurement hardware is different. With a slow clock, aligned
CW-Lite capture, one-hot probe kernels, and a retained binary convolution leak
bank, the recovered images reached 98.41\% bit accuracy and 100\% recognition
accuracy on held-out MNIST traces.

\begin{thebibliography}{9}

\bibitem{wei2018iknow}
L. Wei, B. Luo, Y. Li, Y. Liu, and Q. Xu,
``I Know What You See: Power Side-Channel Attack on Convolutional Neural
Network Accelerators,''
in \emph{ACSAC}, 2018.

\bibitem{cwdocs}
NewAE Technology Inc.,
``ChipWhisperer Documentation,''
\url{https://chipwhisperer.readthedocs.io/}.

\bibitem{mnist}
Y. LeCun, C. Cortes, and C. J. C. Burges,
``The MNIST Database of Handwritten Digits,''
\url{http://yann.lecun.com/exdb/mnist/}.

\end{thebibliography}

\end{document}
